\scanpage{83}
% manuscript line 1
\mspage{- 122 -}n\textsuperscript{o} 316

% manuscript line 2
$V_{(L)}$ est muni canoniquement d'une structure d'espace vectoriel topologique métrisable et complet (\uline{Esp}. \uline{vect}. \uline{top}. chap.I, § 2, n\textsuperscript{o}3, th.2).
% manuscript line 4
Soit $(x_p)_{p\handwritten{\in}z}$ une famille d'éléments de $V$ tels que $x_p=0$ pour $p$ inférieur à un certain entier rationnel. Alors la série $\mathop{\handwritten{\sum}}\nolimits_{\handwritten{-\infty}}^{\handwritten{+\infty}}x_p T^p$ converge dans $V_{(L)}$, et tout élément de $V_{(L)}$ se met sous cette forme de manière unique. (Prendre une base dans $V$).

% manuscript line 8
b) L'espace vectoriel $\handwritten{\unsure{\mathscr L|\mathcal L}}_L(V_{(L)})=(\handwritten{\mathscr L}_K(V))_{(L)}$ est muni canoniquement d'une structure d'espace vectoriel topologique métrisable et complet. Tout endomorphisme de $V_{(L)}$ se met d'une manière unique sous la forme $\mathop{\handwritten{\sum}}\nolimits_{\handwritten{-\infty}}^{\handwritten{+\infty}}u_pT^p$, où $u_p\handwritten{\in}\handwritten{\mathscr L}_K(V)$, et où les $u_p$ sont nuls pour $p$ inférieur à un certain entier rationnel.

% manuscript line 13
c) Si $K$ est de caractéristique $0$, et si $u\handwritten{\in}\handwritten{\mathscr L}_K(V)$, on pose $e^{Tu}={}=1+\frac{1}{1!}\,u\,T+\frac{1}{2!}\,u^2\,T^2+\ldots\handwritten{\in}\handwritten{\mathscr L}_L(V_{(L)})$. Si $x\handwritten{\in}V$ est tel que $u\,x=0$, on a $e^{Tu}.x=x$.

% manuscript line 16
d) Si $W$ est un autre espace vectoriel de dimension finie sur $K$, et si $v\handwritten{\in}\handwritten{\mathscr L}_K(W)$, on a $e^{T(u\handwritten{\otimes}1+1\handwritten{\otimes}v)}=e^{Tu}\handwritten{\otimes}e^{Tv}$.

% manuscript line 18
e) Déduire de c) et d) que, si $V$ est muni d'une structure d'algèbre
% archive sheet number: 82
